\documentclass[11pt,letterpaper]{article}
\usepackage[margin=1in]{geometry}
\usepackage{iftex}
\ifxetex
  \usepackage{fontspec}
  \IfFontExistsTF{Times New Roman}
    {\setmainfont{Times New Roman}}
    {\setmainfont{texgyretermes}[Extension=.otf, UprightFont=*-regular, BoldFont=*-bold,
      ItalicFont=*-italic, BoldItalicFont=*-bolditalic]}
  \setmonofont{DejaVuSansMono}[Extension=.ttf, BoldFont=*-Bold, Scale=0.8]
\else
  \usepackage[T1]{fontenc}
  \usepackage[utf8]{inputenc}
  \usepackage{newtxtext,newtxmath}
  \usepackage[scaled=0.8]{DejaVuSansMono}
\fi
\usepackage{amsmath}
\usepackage{array}
\usepackage{enumitem}
\usepackage{fancyvrb}
\usepackage{xcolor}
\usepackage[hidelinks]{hyperref}
\setlength{\parindent}{0pt}
\setlength{\parskip}{4pt}
\pagestyle{empty}

\usepackage{needspace}
\newcommand{\question}[1]{\needspace{5cm}\vspace{10pt}{\large\bfseries #1}\par\vspace{2pt}}
\newcommand{\blank}{\rule{2.2em}{0.4pt}}
\renewcommand{\arraystretch}{2.2}
\DefineVerbatimEnvironment{code}{Verbatim}{fontsize=\small,frame=none,xleftmargin=1em}

\begin{document}

LING 430 Fall 2026

HW5

Due Thursday Oct 29, 11:59 pm, submit on Canvas

Submission guidelines:
\begin{itemize}[nosep,leftmargin=1.5em]
  \item This is a written homework: work on paper or type it up, whichever you prefer.
  \item Submit one PDF on Canvas under Assignments $>$ HW5; a clear scan or photo of handwriting is fine.
  \item Show your counts and fractions, not only the final numbers.
\end{itemize}

\question{AI usage}

If applicable, list any AI tools you used for this homework, what you asked them, and where the result went into your work. If you used none, write ``none''.

 Declaration:\vspace{1cm}

\question{An HMM tagger from a small corpus}

In this homework you will build an HMM tagger based on a small corpus made up of the following seven sentences. Each sentence has been tagged by hand with two tags, N for noun and V for verb.

\begin{code}
dogs/N bark/V
cats/N watch/V dogs/N
people/N watch/V birds/N
birds/N fly/V
watch/N dogs/N bark/V
dogs/N fly/V
watch/V birds/N
\end{code}

You will work through the tagger on paper in questions 1 to 5: first the three tables, then the tagging of a new sentence. In question 6 you will compare your tagger with the one from nltk. Write every probability as a fraction, for example $6/7$.

\question{1. The start table $\pi$}

Count how many sentences begin with each tag and fill in $\pi$, the probability that a sentence begins with N or with V.

\begin{tabular}{|l|>{\centering\arraybackslash}p{2.2cm}|}
\hline
$\pi$ & \\ \hline
N &  \\ \hline
V &  \\ \hline
\end{tabular}

\par\vspace{6pt}

\question{2. The transition table A}

Count the tag bigrams and fill in A, the probability of each tag given the previous tag. The rows are the previous tag. The last tag of a sentence does not begin a bigram.

\begin{tabular}{|l|>{\centering\arraybackslash}p{2.2cm}|>{\centering\arraybackslash}p{2.2cm}|}
\hline
A & N & V \\ \hline
N &  &  \\ \hline
V &  &  \\ \hline
\end{tabular}

\par\vspace{6pt}

\question{3. The emission table B}

Count the words under each tag and fill in B, the probability of each word given its tag.

\begin{tabular}{|l|*{7}{>{\centering\arraybackslash}p{1.35cm}|}}
\hline
B & dogs & cats & people & birds & watch & bark & fly \\ \hline
N &  &  &  &  &  &  &  \\ \hline
V &  &  &  &  &  &  &  \\ \hline
\end{tabular}

\par\vspace{6pt}

\question{4. Tagging a new sentence}

The sentence \emph{watch dogs bark} has two readings: \emph{watchdogs bark}, N N V, and \emph{watch the dogs bark}, V N V. Fill in the six factors of each path from your tables and multiply them to a decimal. Which reading does your tagger prefer, and by how much?

\begin{align*}
\text{N N V:}\quad & P(\text{N}) \times P(\text{watch} \mid \text{N}) \times P(\text{N} \mid \text{N}) \times P(\text{dogs} \mid \text{N}) \times P(\text{V} \mid \text{N}) \times P(\text{bark} \mid \text{V}) \\
&= \blank \times \blank \times \blank \times \blank \times \blank \times \blank = \blank \\[6pt]
\text{V N V:}\quad & P(\text{V}) \times P(\text{watch} \mid \text{V}) \times P(\text{N} \mid \text{V}) \times P(\text{dogs} \mid \text{N}) \times P(\text{V} \mid \text{N}) \times P(\text{bark} \mid \text{V}) \\
&= \blank \times \blank \times \blank \times \blank \times \blank \times \blank = \blank
\end{align*}
\par\vspace{4pt}Answer:\par\vspace{1cm}

\question{5. The other paths}

With two tags and three words there are eight possible tag sequences for \emph{watch dogs bark}. Explain why the other six have probability 0.

\par\vspace{4pt}Answer:\par\vspace{2cm}

\question{6. Comparing with nltk}

Paste the following into a code cell in Colab and run it. It tags \emph{watch dogs bark} and the seven sentences of the corpus with the nltk tagger from class 9a.

\begin{code}
import nltk
nltk.download("averaged_perceptron_tagger_eng", quiet=True)

sentences = ["watch dogs bark",
             "dogs bark", "cats watch dogs", "people watch birds", "birds fly",
             "watch dogs bark", "dogs fly", "watch birds"]
for s in sentences:
    print(nltk.pos_tag(s.split()))
\end{code}

Treat NN, NNS and NNP as N, and VB, VBP and VBZ as V. Does nltk choose the same reading of \emph{watch dogs bark} as your tagger? Where does it disagree with the hand tags in the corpus? Then reflect in three to five sentences: the nltk tagger was trained on a million words of newspaper text and yours on seven sentences. Which one would you trust on this sentence, and what would it take to change your mind?

\par\vspace{4pt}Answer:\par\vspace{4cm}

\end{document}
